\scanpage{335}
\mspage{74}
% archive sheet number: 334
% ms. 1–2
\deleted{flèches sont les flèches \uline{inversibles} de $Y$. Donc $f$ induit}
% ms. 3
\[\deleted{f_{\unsure{0|o}}:X\longrightarrow{}}\]
% ms. 4–5
des flèches inversibles de $X'$ (lequel ensemble satisfait bien aux deux conditions d'admissibilité).
% ms. 6–7
Donc \deleted{\unsure{par}} ce qui précède, on a un diagramme commutatif
% ms. 8–10
\diagram{\begin{tikzcd}[ampersand replacement=\&,row sep=large,column sep=large]
X\arrow[r,"k"]\arrow[d,"f"'] \& \Pi_\Sigma(X)\arrow[dl,dashed,"\tilde f"]\arrow[d,"{\bar f=\deleted{\unsure{f}}\Pi_\Sigma(f)}"]\\
X'\arrow[r,"{k'}","\sim" description] \& \Pi_{\Sigma'}(X'){,}
\end{tikzcd}}
% ms. 11–15
où $\bar f$ est \unsure{uniquement} déterminé. Comme $k'$ est un isomorphisme (comme on a vu tantôt, puisque $\Sigma'$ \unsure{formé|formée} de fl. \unsure{inv.|inver.}), l'existence et l'unicité de $\deleted{\unsure{\Pi}}\unsure{\tilde f|\bar f}$ \unsure{équivalent|équivaut} à la conclusion annoncée, cqfd.
% ms. between 15 and 16
\figure{\begin{tikzpicture}[x=1cm,y=1cm]\draw (0,.32) -- (.23,0);\end{tikzpicture}}

% ms. 16–19
\uline{Encore un \unsure{sorite}} sur les classes d'homotopie d'éléments d'une catégorie $X$. De façon générale, on voit que \deleted{les \unsure{cl}} $X$ et $\Pi_\Sigma(X)=X\Sigma^{-1}$ ont \unsure{$\Longrightarrow$|mêmes} \deleted{\unsure{classes de}} composantes connexes, en \unsure{ce|si}
